% SPDX-FileCopyrightText: 2016-2026 Lari Natri <lari.natri@iki.fi>
% SPDX-License-Identifier: GPL-3.0-or-later
%
% Book display and text defaults.
% This file is part of the 'ulsbs' package.
%
% Loaded after preliminary language selection and ulsbs-utilities.code.tex.
% Requires the boolean factory and ulIfLanguage; owns display switches,
% the initial verse-indent dimension and book/metadata strings.
% The 1em indent is captured before typography selects the document fonts.

\ProvidesFile{ulsbs-book-defaults.code.tex}[Book display and text defaults]

% === Display switches ===

% Introduce toggles for some global features and set their defaults. If you
% want a specific feature enabled or disabled, set its public setter to
% true or false in the main document preamble (after including this package).
% Example: \ulShowTags{false} would disable showing tags in song preludes, and
%          \ulShowTags{true} would enable showing them.
% Query with the real predicate \ifulShowTags ...\else ...\fi. Setters obey
% ordinary grouping; do not predefine or renew these commands/conditionals.

% Toggle for showing extra info in song preludes. Set the song metadata with
% \beginsong{Title}[ex={this song is good}].
\ulsbs@newbool\ifulShowExtraInfo\ulsbs@showextrue\ulsbs@showexfalse\ulShowExtraInfo\ulsbs@showextrue

% Toggle for showing the dedication in song preludes. Set the song metadata with
% \beginsong{Title}[off={My Muse}].
\ulsbs@newbool\ifulShowOfferedTo\ulsbs@showofftrue\ulsbs@showofffalse\ulShowOfferedTo\ulsbs@showofftrue

% Toggle for showing key info in song preludes. To set it for a song,
% use \beginsong{Title}[key={Am}]. Make sure that
% it is correct, that the key set here corresponds to the actual key displayed
% in the final document (after possible transpose).
\ulsbs@newbool\ifulShowKey\ulsbs@showkeytrue\ulsbs@showkeyfalse\ulShowKey\ulsbs@showkeytrue

% Toggle for showing good keys in song preludes. To set it/them for a song,
% use \beginsong{Title}[gk={Am, Gm--Bm}]. The first
% is the preferred one, and then comes the range of keys suitable for singing
% in such a way that the melody falls within (including) a G note in the low
% end, and a B note of the next octave in the high end. If such a range is
% impossible, put the keys in the range in parenthesis.
\ulsbs@newbool\ifulShowGoodKeys\ulsbs@showgoodkeystrue\ulsbs@showgoodkeysfalse\ulShowGoodKeys\ulsbs@showgoodkeystrue

% Toggle for showing phase info in song preludes. To set it for a song,
% use \beginsong{Title}[ph={I, II}].
\ulsbs@newbool\ifulShowPhase\ulsbs@showphtrue\ulsbs@showphfalse\ulShowPhase\ulsbs@showphtrue

% Toggle for showing phase info in the TOC:
\ulsbs@newbool\ifulShowPhaseInToc\ulsbs@showphintoctrue\ulsbs@showphintocfalse\ulShowPhaseInToc\ulsbs@showphintoctrue

% Toggle for showing tags in song preludes. To set them for a song,
% use \beginsong{Title}[tags={love, light}].
% Consumer books may define their tag vocabulary in include/tags.txt.
\ulsbs@newbool\ifulShowTags\ulsbs@showtagstrue\ulsbs@showtagsfalse\ulShowTags\ulsbs@showtagstrue

% Toggle for showing melody notes in songs. Write pitches inside a melody
% block in a chord definition, e.g. \[<C D>Am]. To set them temporarily on or
% off within a (part of a) song, use macros \notesoff and \noteson.
\ulsbs@newbool\ifulShowNotes\ulsbs@shownotestrue\ulsbs@shownotesfalse\ulShowNotes\ulsbs@shownotestrue

% Toggle for showing beat marks in songs. Write * inside a melody block in
% a chord definition, e.g. \[<* C>Am]. To set them temporarily on or
% off within a (part of a) song, use macros \beatsoff and \beatson.
\ulsbs@newbool\ifulShowBeats\ulsbs@showbeatstrue\ulsbs@showbeatsfalse\ulShowBeats\ulsbs@showbeatstrue

% Toggle for showing authors in the song preludes. Personal opinion: the
% authors should definitely be shown in all cases. This option is here for
% the sake of completeness. :)
\ulsbs@newbool\ifulShowAuthors\ulsbs@showauthtrue\ulsbs@showauthfalse\ulShowAuthors\ulsbs@showauthtrue

% Toggle for showing Lilypond music scores
\ulsbs@newbool\ifulShowLilypond\ulsbs@showlilypondtrue\ulsbs@showlilypondfalse\ulShowLilypond\ulsbs@showlilypondtrue

% Toggle for showing passages
\ulsbs@newbool\ifulShowPassage\ulsbs@showpassagetrue\ulsbs@showpassagefalse\ulShowPassage\ulsbs@showpassagetrue

% Toggle for showing explanations
\ulsbs@newbool\ifulShowExplanation\ulsbs@showexplanationtrue\ulsbs@showexplanationfalse\ulShowExplanation\ulsbs@showexplanationtrue

% Toggle for showing translations
\ulsbs@newbool\ifulShowTranslation\ulsbs@showtranslationtrue\ulsbs@showtranslationfalse\ulShowTranslation\ulsbs@showtranslationtrue

% Toggle for showing feelers
\ulsbs@newbool\ifulShowFeeler\ulsbs@showfeelertrue\ulsbs@showfeelerfalse\ulShowFeeler\ulsbs@showfeelertrue

% Toggle showing of alternate chord sets
\ulsbs@newbool\ifulShowAltChords\ulsbs@showaltchordstrue\ulsbs@showaltchordsfalse\ulShowAltChords\ulsbs@showaltchordstrue

% Toggle transposing alternate chord sets with the current \transpose factor.
% This is disabled by default because alternate sets usually represent chords
% in a different key and should remain unchanged when the main chords move.
\ulsbs@newbool\ifulTransposeAltChords\ulsbs@transposealtchordstrue\ulsbs@transposealtchordsfalse\ulTransposeAltChords\ulsbs@transposealtchordsfalse

% Toggle appending the (optional) chapter image or symbol appended to the
% chapter name wherever it shows up, e.g. in TOC and page headers. The image is
% never shown in PDF index though, as only characters are supported there.
\ulsbs@newbool\ifulShowImageInChapterName\ulsbs@showimageinchapternametrue\ulsbs@showimageinchapternamefalse\ulShowImageInChapterName\ulsbs@showimageinchapternametrue

% Toggle showing the (optional) chapter image or symbol on the front page of
% a chapter.
\ulsbs@newbool\ifulShowImageOnChapterFrontPage\ulsbs@showimageonchapterfrontpagetrue\ulsbs@showimageonchapterfrontpagefalse\ulShowImageOnChapterFrontPage\ulsbs@showimageonchapterfrontpagetrue

% Toggle uppercasing the long chapter title shown on chapter front pages
\ulsbs@newbool\ifulUppercaseLongChapterTitle\ulsbs@upcaselongchaptertitletrue\ulsbs@upcaselongchaptertitlefalse\ulUppercaseLongChapterTitle\ulsbs@upcaselongchaptertitlefalse

% Toggle uppercasing the short chapter title shown in TOC and page headers
\ulsbs@newbool\ifulUppercaseShortChapterTitle\ulsbs@upcaseshortchaptertitletrue\ulsbs@upcaseshortchaptertitlefalse\ulUppercaseShortChapterTitle\ulsbs@upcaseshortchaptertitletrue

% Toggle uppercasing the section title in page headers
\ulsbs@newbool\ifulUppercaseSectionTitleInHeader\ulsbs@upcasesectiontitleinheadertrue\ulsbs@upcasesectiontitleinheaderfalse\ulUppercaseSectionTitleInHeader\ulsbs@upcasesectiontitleinheaderfalse

% Toggle uppercasing the book name in page headers
\ulsbs@newbool\ifulUppercaseBookTitleInHeader\ulsbs@upcasebooktitleinheadertrue\ulsbs@upcasebooktitleinheaderfalse\ulUppercaseBookTitleInHeader\ulsbs@upcasebooktitleinheadertrue

% Toggle use of chapter colors in page corners
\ulsbs@newbool\ifulUseChapterColors\ulsbs@usechaptercolorstrue\ulsbs@usechaptercolorsfalse\ulUseChapterColors\ulsbs@usechaptercolorstrue

% Internal current-style switch.  \ulMelodyStyleMapping resolves semantic roles
% to these two existing configurable visual styles.
\newif\ifulsbs@altmnstyle\ulsbs@altmnstylefalse

% If true, some characters on chord line are adjusted: currently parenthesis
% are scaled down and moved slightly upwards to avoid collision with lyrics.
% This does not affect chords displayed within \altchords{} macro.
\ulsbs@newbool\ifulAdjustChordChars\ulsbs@adjustchordcharstrue\ulsbs@adjustchordcharsfalse\ulAdjustChordChars\ulsbs@adjustchordcharstrue

% === Verse indentation and song-end mark ===

% Indentation that is optionally added to verses, if started with
% \beginverse[n], where n is a plain factor number that is multiplied with
% the \ulVerseIndentWidth to determine the optional extra indent depth of the
% verse in question.
% This 1em is captured now, before document typography. Override the allocated
% dimension after loading with \setlength{\ulVerseIndentWidth}{<dimension>}.
\newdimen\ulVerseIndentWidth\ulVerseIndentWidth=1em\relax

\providecommand{\ulSongEndMark}{}

% === Book strings and PDF metadata defaults ===

% Set defaults for book titles if they're not yet defined. One should define
% them with \providecommand before loading or \renewcommand afterward for
% front matter and headers. PDF metadata is consumed by the shared style's
% load-time \hypersetup, so define metadata/title/version before loading (or
% explicitly update \hypersetup later). The supported macros are:
%   \ulMainBookTitle
%   \ulSubBookTitle
%   \ulSubSubBookTitle
%   \ulThisBookTitle
%   \ulBookMotto
%   \ulBookByText
%   \ulBookVersionPreText
%   \ulBookVariantPreText
%   \ulBookDefaultVariantText
%   \ulBookLyricVariantText
%   \ulBookCharangoVariantText
%   \ulBookBassclefVariantText
%   \ulBookWebsiteLink
%   \ulBookWebsiteLinkQrImage (this is image filename in content/img without path)
%   \ulBookCompiledByText
%   \ulImprintPageFootnote
% All are optional; the main title also has a fallback below.

% Set defaults for PDF metadata, if they are not defined. One should define
% them in the main file before including this package.
\providecommand*{\ulBookAuthor}{}
\providecommand*{\ulBookSubject}{Songbook}
\providecommand*{\ulBookKeywords}{songbook, songs, ulsbs}

\providecommand{\ulBookMotto}{}
\providecommand{\ulSubBookTitle}{}
\providecommand{\ulSubSubBookTitle}{}
\providecommand{\ulBookWebsiteLink}{}
\providecommand{\ulBookWebsiteLinkQrImage}{}
\providecommand{\ulBookCompiledByText}{\ifdefempty{\ulBookAuthor}{}{Compiled by: \ulBookAuthor}}
\providecommand{\ulImprintPageFootnote}{}

% To have the portuguese version of these default strings, one must define
% \newcommand{\ulPreliminaryLanguage}{portuguese} before including this file.
\ulIfLanguage{portuguese}{%
  \providecommand{\ulMainBookTitle}{Songbook Sem Nome}
  \providecommand{\ulBookByText}{Por:\par\vspace{0.618034ex}\textbf{\Large humanidade}}
  \providecommand{\ulBookVersionPreText}{Versão:}
  \providecommand{\ulBookVariantPreText}{Variante:}
  \providecommand{\ulBookDefaultVariantText}{padrão}
  \providecommand{\ulBookLyricVariantText}{somente letras}
  \providecommand{\ulBookCharangoVariantText}{por charango}
  \providecommand{\ulBookBassclefVariantText}{com clave de fá}
}{%
  \providecommand{\ulMainBookTitle}{Songbook Without Name}
  \providecommand{\ulBookByText}{By:\par\vspace{0.618034ex}\textbf{\Large humankind}}
  \providecommand{\ulBookVersionPreText}{Version:}
  \providecommand{\ulBookVariantPreText}{Variant:}
  \providecommand{\ulBookDefaultVariantText}{default}
  \providecommand{\ulBookLyricVariantText}{lyrics only}
  \providecommand{\ulBookCharangoVariantText}{for charango}
  \providecommand{\ulBookBassclefVariantText}{with bass clef}
}

% The book title to display on top of each verso page (if doesn't exist, it is
% constructed here from \ulMainBookTitle and optional \ulSubBookTitle):
\providecommand{\ulThisBookTitle}{% main title (+ subtitle, if defined)
  \ulMainBookTitle\ifdefempty{\ulSubBookTitle}{}{{: \ulSubBookTitle}}%
}
% The default is the current date
\providecommand{\ulBookVersion}{%
  %2024-08-12% for release: change and uncomment this line, and comment the next
  \the\year-\ifnum\month<10 0\fi\the\month-\ifnum\day<10 0\fi\the\day%
}
